% =====================================================================
%  Section 3: Method — Trajectory-Anchored Dynamic Selection (TADS)
% =====================================================================
\section{Method: Trajectory-Anchored Dynamic Selection}
\label{sec:method}

This section defines the data-selection problem
(Sec.~\ref{sec:problem-setup}), explains why multi-task selection
benefits from a structural anchor (Sec.~\ref{sec:method-motivation}),
presents the \textsc{TADS} selection method
(Sec.~\ref{sec:method-core}), and states the anchor-stability
guarantee (Sec.~\ref{sec:method-theory}).

\subsection{Problem setup}
\label{sec:problem-setup}

Let $\mathcal{D} = \{x_i\}_{i=1}^N$ be a candidate instruction-tuning
pool, where each $x_i$ is an instruction--response pair tokenised
as a single sequence, and let $f_\theta$ be a base LLM with
parameters $\theta$. Given a selection budget $B \le N$
(equivalently, a selection ratio $\rho = B/N$ with $0 < \rho \le 1$),
the goal is to choose a subset $\mathcal{S} \subset \mathcal{D}$
with $|\mathcal{S}| = B$ such that fine-tuning $f_\theta$ on
$\mathcal{S}$ maximises performance on a set of evaluation tasks
$\{T_1, \ldots, T_M\}$ spanning multiple capabilities---knowledge,
reasoning, code, and multilingual understanding in our experiments.
We work in the \emph{multi-task} setting: a single subset
$\mathcal{S}$ is shared across all $M$ tasks rather than re-selected
per task, reflecting the practical regime in which a model is
fine-tuned once and deployed across many downstream uses.

\paragraph{Token-level state dynamics.}
For a hidden layer $l \in \{1, \ldots, L\}$ of $f_\theta$ and an
input sequence $x$ with $K_x$ valid (non-padding) tokens, we write
$h_l^{(k)}(x;\theta) \in \mathbb{R}^d$ for the layer-$l$ hidden
state at the $k$-th token position, where $d$ is the model's hidden
dimension. A causal language model processes $x$ as a
token-by-token trajectory of hidden states; the transition between
consecutive positions is
\begin{equation}
  s_l^{(k)}(x;\theta)
  \;:=\;
  h_l^{(k)}(x;\theta) - h_l^{(k-1)}(x;\theta),
  \qquad k = 2, \ldots, K_x .
  \label{eq:transition}
\end{equation}
Recent work~\citep{tang2025sde} shows that these hidden-state
transitions can carry reasoning information not preserved by
surface tokens alone. For per-sample selection, however, each
candidate must be reduced to a comparable scalar score, so we
exploit two complementary sample-level aggregations of this
trajectory.

\paragraph{Contextualisation delta.}
Write $h_l^{\mathrm{first}}(x;\theta) := h_l^{(1)}(x;\theta)$ and
$h_l^{\mathrm{last}}(x;\theta) := h_l^{(K_x)}(x;\theta)$ for the
first and last valid-token hidden states.\footnote{The first and
last valid positions are the sequence-start and end markers; what
matters is the hidden state \emph{at} those positions---before any
content is processed, and after the full causal pass over
$x$---not the markers themselves.} The \emph{contextualisation
delta} of $x$ at layer $l$ is
\begin{equation}
  \Delta h_l(x;\theta) \;:=\;
  h_l^{\mathrm{last}}(x;\theta) - h_l^{\mathrm{first}}(x;\theta)
  \;=\; \sum_{k=2}^{K_x} s_l^{(k)}(x;\theta),
  \label{eq:delta}
\end{equation}
the telescoping sum of token-level state transitions
(Eq.~\ref{eq:transition}). The delta isolates the cumulative effect
of processing $x$ on the layer's representation---``what the sample
did'' to the model---independent of any absolute reference point.
This is the same per-sample quantity that capability-direction
methods such as NAIT~\citep{nait2026} extract from in-domain seed
samples; the distinction in \textsc{TADS} lies in the \emph{source}
of $\{\Delta h_l\}$---the candidate pool itself, at each training
epoch, rather than a curated external seed.

\paragraph{Sequence-mean activation.}
The \emph{sequence-mean activation} of $x$ at layer $l$ is
\begin{equation}
  \bar h_l(x;\theta) \;:=\;
  \frac{1}{K_x}\sum_{k=1}^{K_x} h_l^{(k)}(x;\theta),
  \label{eq:meanact}
\end{equation}
the average representation across all valid token positions---a
padding-invariant aggregate that does not depend on any single
endpoint and smooths over local token-level idiosyncrasies.

\subsection{Why multi-task selection benefits from a structural anchor}
\label{sec:method-motivation}

RL-based dynamic selectors following the learned advantage estimator
recipe~\cite{schulman2017ppo, schulman2016gae} compose two
per-candidate signals at refresh step $t$---a difficulty term
$\mathcal{L}_i^{(t)}$ (response-token cross-entropy of the current
model on $x_i$) and an uncertainty term $\mathcal{H}_i^{(t)}$ (mean
predictive entropy over response tokens)---into a single composite
reward
\begin{equation}
  \mathcal{R}_i^{(t)} \;=\;
  w^{(t)} \, \mathcal{L}_i^{(t)}
  + \bigl(1 - w^{(t)}\bigr)\, \mathcal{H}_i^{(t)},
  \label{eq:composite}
\end{equation}
with a variance-ratio weight
\begin{equation}
  w^{(t)} \;=\;
  \frac{\mathrm{Var}_{x \in \mathcal{D}}(\mathcal{L}^{(t)})}
       {\mathrm{Var}_{x \in \mathcal{D}}(\mathcal{L}^{(t)}) +
        \mathrm{Var}_{x \in \mathcal{D}}(\mathcal{H}^{(t)}) + \varepsilon},
  \label{eq:w}
\end{equation}
where each variance is taken over the entire candidate pool
$\mathcal{D}$ at the current epoch. We adopt this construction
unchanged as the composite-reward component of our dynamic selection
pipeline; the dataset-level definition of $w^{(t)}$ is discussed in
Sec.~\ref{sec:method-core}.

\paragraph{The multi-task variance trap.}
By the law of total variance, decomposing $\mathcal{L}^{(t)}$ over the
$M$ tasks that compose $\mathcal{D}$,
\begin{equation}
  \mathrm{Var}_{x \in \mathcal{D}}(\mathcal{L}^{(t)})
  \;=\;
  \underbrace{\mathbb{E}_{m}\!\bigl[\mathrm{Var}_{x \in T_m}(\mathcal{L}^{(t)})\bigr]}_{\text{within-task}}
  \;+\;
  \underbrace{\mathrm{Var}_{m}\!\bigl[\mathbb{E}_{x \in T_m}(\mathcal{L}^{(t)})\bigr]}_{\text{between-task}},
  \label{eq:total-variance}
\end{equation}
and analogously for $\mathcal{H}^{(t)}$. In a single-task pool the
between-task term vanishes, leaving $w^{(t)}$ a stable summary of
within-task informativeness. In a multi-task pool with capabilities
of disparate difficulty, the between-task term \emph{dominates}:
it depends on the task proportions in $\mathcal{D}$ and on the
relative scales of $\mathcal{L}^{(t)}$ and $\mathcal{H}^{(t)}$ across
capabilities. Consequently $w^{(t)}$ becomes a sensitive function of
the task-mix, drifting between runs at fixed hyperparameters, and the
scalar reward $\mathcal{R}_i^{(t)}$ can no longer separate samples by
\emph{which capability} they exercise.

\paragraph{A bounded, pool-level signal absorbs the drift.}
The trajectory anchor $v_l^{(t)}$ (defined in
Sec.~\ref{sec:method-core}) is a pool-level summary: the top
principal direction of the contextualisation deltas across a probe
subset, not a per-sample statistic. Two properties make it a
remedy for the drift above. \emph{Geometrically}, $v_l^{(t)}$
captures the capability subspace the model is most actively shaping
at this epoch and layer---a property of the parameter--data
geometry, independent of how the task-mix balances the two terms in
Eq.~\ref{eq:total-variance}. \emph{Analytically}, its
epoch-to-epoch change is uniformly bounded by the learning rate,
$\|v_l^{(t+1)} - v_l^{(t)}\| \le (2 C_\Sigma/\gamma_{\min})\,\eta_t$
(Theorem~\ref{thm:stability}), whereas $w^{(t)}$ has no such bound.
The \textsc{TADS} score (Sec.~\ref{sec:method-core}) therefore
combines the task-mix-drifting base signal
$\mathcal{R}_i^{(t)} \cdot a_i^{(t)}$ with this
stability-controlled structural signal, each at the variance scale
where it is informative.

\begin{figure}[t]
\centering
\begin{tikzpicture}[scale=0.95, every node/.style={font=\small}]
  \draw[gray!30, very thin, -{Latex[length=2mm]}] (-0.1,0) -- (6.4,0);
  \draw[gray!30, very thin, -{Latex[length=2mm]}] (0,-0.1) -- (0,3.8);
  \node[gray!50, font=\scriptsize] at (6.5,-0.18) {$\mathbb{R}^d$};
  \filldraw[black] (0,0) circle (1.2pt);
  \foreach \x/\y in
    {3.8/1.2, 4.2/1.6, 3.5/1.8, 4.5/2.0, 3.0/1.3, 4.0/2.3, 3.7/1.5, 4.6/1.4, 3.3/1.6}
  {
    \draw[-{Latex[length=1.6mm]}, orange!60, thin, opacity=0.75]
      (0,0) -- (\x,\y);
  }
  \node[orange!80!black, font=\footnotesize, align=left, anchor=north west] at (3.5,0.85)
    {Probe deltas\\
     $\{\Delta h_l(x;\theta_{t-1})\}_{x\in\widetilde{\mathcal{D}}_t}$};
  \draw[-{Latex[length=2.6mm]}, blue!75!black, ultra thick]
    (0,0) -- (5.0,2.0);
  \node[blue!75!black, font=\bfseries] at (5.35,2.25) {$v_l^{(t)}$};
  \draw[-{Latex[length=2.2mm]}, blue!35, dashed, semithick]
    (0,0) -- (5.0,1.65);
  \node[blue!50, font=\scriptsize] at (5.40,1.45) {$v_l^{(t-1)}$};
  \node[gray!60!black, font=\scriptsize, align=left] at (6.4,1.85)
    {Thm.~\ref{thm:stability}:\\[-1pt]
     $\le \tfrac{2C_\Sigma}{\gamma_{\min}}\,\eta_t$};
  \filldraw[red!75!black] (2.2,2.85) circle (2pt);
  \node[red!75!black, font=\footnotesize] at (1.45,3.05)
    {$\bar h_l(x_i;\theta_{t-1})$};
  \draw[red!75!black, dotted, thick] (2.2,2.85) -- (2.88,1.15);
  \filldraw[red!75!black] (2.88,1.15) circle (1.5pt);
  \node[red!75!black, font=\scriptsize, anchor=west]
    at (1.7, 1.2) {$\mathrm{align}_i^{(t)}$};
\end{tikzpicture}
\caption{\textbf{Trajectory anchor at epoch~$t$.}
Orange arrows are the contextualisation deltas
$\Delta h_l(x;\theta_{t-1})$ from a probe subset
$\widetilde{\mathcal{D}}_t$---the \emph{sequence-level} trajectory
inside each sample. The bold blue arrow $v_l^{(t)}$ is their top
principal direction (the anchor); the dashed $v_l^{(t-1)}$ is the
previous epoch's anchor, the \emph{epoch-level} trajectory whose
deviation is bounded by Theorem~\ref{thm:stability}. A candidate's
alignment is the projection of its sequence-mean activation
$\bar h_l(x_i;\theta_{t-1})$ onto $v_l^{(t)}$.}
\label{fig:anchor}
\end{figure}

\subsection{Trajectory-Anchored Selection}
\label{sec:method-core}

\textsc{TADS} scores each candidate at epoch~$t$ by reweighting a
base dynamic ranking rule with a structural anchor extracted from
the model's own training trajectory. We describe anchor extraction,
candidate alignment, the selection score, and the full algorithm in
turn.

\paragraph{Probe subset and trajectory anchor.}
At the start of epoch~$t$ we draw a probe subset
$\widetilde{\mathcal{D}}_t \subset \mathcal{D}$ uniformly at random
with $|\widetilde{\mathcal{D}}_t| = 1{,}024$, re-sampled each epoch
so that no fixed subset biases the anchor. For each layer $l$ we
collect the contextualisation deltas
$\{\Delta h_l(x;\theta_{t-1}) : x \in \widetilde{\mathcal{D}}_t\}$
(Eq.~\ref{eq:delta}) and form their mean and centred covariance:
\begin{align}
  \bar{\Delta h}_l^{(t)}
   &:= \frac{1}{|\widetilde{\mathcal{D}}_t|}
      \sum_{x \in \widetilde{\mathcal{D}}_t} \Delta h_l(x;\theta_{t-1}),
   \label{eq:mu} \\[2pt]
  \Sigma_l^{(t)}
   &:= \frac{1}{|\widetilde{\mathcal{D}}_t|}
      \sum_{x \in \widetilde{\mathcal{D}}_t}
      \bigl(\Delta h_l(x;\theta_{t-1}) - \bar{\Delta h}_l^{(t)}\bigr)
      \bigl(\Delta h_l(x;\theta_{t-1}) - \bar{\Delta h}_l^{(t)}\bigr)^{\!\top}.
   \label{eq:sigma}
\end{align}
The trajectory anchor for layer $l$ is the top unit eigenvector of
$\Sigma_l^{(t)}$,
\begin{equation}
  v_l^{(t)} \;=\;
  \arg\max_{\|v\| = 1}\; v^{\!\top} \Sigma_l^{(t)}\, v,
  \label{eq:anchor-def}
\end{equation}
sign-calibrated so that
$\langle v_l^{(t)},\, \bar{\Delta h}_l^{(t)} \rangle > 0$, which
removes the PCA $\pm v$ ambiguity.\footnote{This sign calibration
also serves as Assumption~3 of Theorem~\ref{thm:stability}; see
Sec.~\ref{sec:method-theory}.} Intuitively, $v_l^{(t)}$ is the
direction that best explains how layer $l$ responds to the probe
samples at the current parameter state---the \emph{capability
direction} the model is most actively shaping at this epoch. It
depends only on $\theta_{t-1}$ and the candidate pool: no in-domain
seed, no labels, no auxiliary model.

\paragraph{Candidate alignment.}
For each candidate $x_i$, the per-layer alignment is the projection
of its sequence-mean activation (Eq.~\ref{eq:meanact}) onto the
anchor, $\langle \bar h_l(x_i;\theta_{t-1}), v_l^{(t)} \rangle$. The
aggregate alignment averages over layers,
\begin{equation}
  \mathrm{align}_i^{(t)}
  \;:=\;
  \frac{1}{L}\sum_{l=1}^{L}
     \bigl\langle \bar h_l(x_i;\theta_{t-1}),\, v_l^{(t)} \bigr\rangle,
  \label{eq:align-raw}
\end{equation}
and is then min--max normalised across the pool,
\begin{equation}
  \widetilde{\mathrm{align}}_i^{(t)}
  \;:=\;
  \frac{\mathrm{align}_i^{(t)} - \min_j \mathrm{align}_j^{(t)}}
       {\max_j \mathrm{align}_j^{(t)} - \min_j \mathrm{align}_j^{(t)}}
  \;\in\; [0, 1].
  \label{eq:align}
\end{equation}
Averaging over all valid positions makes the score
padding-invariant; the layer subset used for alignment is one of
the design choices examined in App.~\ref{app:inprogress}.

\paragraph{Selection score.}
The advantage-estimator pipeline provides two time-dependent signals for
each candidate at refresh step $t$: the composite reward
$\mathcal{R}_i^{(t)}$ (Eq.~\ref{eq:composite}) and a advantage estimate
$a_i^{(t)} \in [0,1]$. We use their product as a fixed
\emph{reward--advantage diagnostic consensus} base ranking: a candidate
receives a high base score only when both the scalar reward signal
and the learned advantage estimate support its selection. \textsc{TADS}
then multiplies this base ranking by the trajectory-anchor factor:
\begin{equation}
\boxed{
\,s_i^{(t)} \;=\;
\underbrace{\mathcal{R}_i^{(t)} \cdot a_i^{(t)}}_{\substack{\text{base factor}\\[1pt] \text{(high-variance, sample-level)}}}
\;\cdot\;
\underbrace{\bigl(1 + \lambda \cdot \widetilde{\mathrm{align}}_i^{(t)}\bigr)}_{\substack{\text{anchor factor} \,\in\, [1,\,1+\lambda]\\[1pt] \text{(low-variance, pool-level)}}}\,
}
\label{eq:score}
\end{equation}
with a single hyperparameter $\lambda \ge 0$. The multiplicative
form carries three properties. \emph{(i) Clean ablation:}
$\lambda{=}0$ makes the anchor factor $\equiv 1$, recovering the
fixed reward--advantage consensus base ranking, so any gain at
$\lambda > 0$ is attributable to the anchor alone. \emph{(ii)
Variance-scale decoupling:} the base factor inherits the task-mix
drift of Eq.~\ref{eq:total-variance} while the anchor factor, a
pool-level PCA statistic, does not---a separation that additive
composition would forfeit. \emph{(iii) Stability transfer:} because
the anchor factor is bounded in $[1,1+\lambda]$ and stable
(Theorem~\ref{thm:stability}), the anchor-induced change in
$s_i^{(t)}$ is bounded by $O(\lambda\eta_t)$; the proof is given in
App.~\ref{app:score-stability}.

\paragraph{Implementation note: dataset-level $w^{(t)}$.}
The variance-ratio weight $w^{(t)}$ (Eq.~\ref{eq:w}) must be computed
at the \emph{dataset level}, over all $N$ candidates. A naive
per-mini-batch implementation recomputes
$\mathrm{Var}(\mathcal{L}^{(t)})$ and
$\mathrm{Var}(\mathcal{H}^{(t)})$ inside each batch; with
$|\mathcal{B}|{=}1$ (a common memory-efficient default) the per-batch
variance is identically zero, so $w^{(t)} \equiv 0$ and the composite
reward silently collapses to
$\mathcal{R}_i^{(t)} = \mathcal{H}_i^{(t)}$ for the entire run.
\textsc{TADS} instead accumulates
$\{\mathcal{L}_i^{(t)}\}$ and $\{\mathcal{H}_i^{(t)}\}$ across all
batches and computes $w^{(t)}$ once per epoch. The per-batch variant
is included among the diagnostic experiments of
App.~\ref{app:inprogress}.

\begin{algorithm}[t]
\caption{Trajectory-Anchored Dynamic Selection (TADS)}
\label{alg:tads}
\begin{algorithmic}[1]
\REQUIRE Pool $\mathcal{D}$, base model $f_{\theta_0}$, budget $B$,
         layers $\{1,\ldots,L\}$, weight $\lambda$, epochs $T$,
         probe size $|\widetilde{\mathcal{D}}|$
\FOR{$t = 1, \dots, T$}
  \STATE Sample probe set
         $\widetilde{\mathcal{D}}_t \subset \mathcal{D}$ uniformly at random
  \FOR{$l = 1, \dots, L$}
    \STATE Compute deltas
           $\{\Delta h_l(x; \theta_{t-1})\}_{x \in \widetilde{\mathcal{D}}_t}$,
           covariance $\Sigma_l^{(t)}$
           \hfill (Eqs.~\ref{eq:delta},~\ref{eq:sigma})
    \STATE $v_l^{(t)} \leftarrow$ top-1 eigenvector of $\Sigma_l^{(t)}$,
           sign-calibrated
  \ENDFOR
  \STATE For all $i$: aggregate
         $\widetilde{\mathrm{align}}_i^{(t)}$
         \hfill (Eq.~\ref{eq:align})
  \STATE Compute composite reward and advantage estimate (App.~B); collect
         $\{\mathcal{L}_i^{(t)}, \mathcal{H}_i^{(t)}, a_i^{(t)}\}_{i=1}^N$
  \STATE Compute pool-level $w^{(t)}$, composite reward
         $\mathcal{R}_i^{(t)}$
         \hfill (Eqs.~\ref{eq:w},~\ref{eq:composite})
  \STATE Score
         $s_i^{(t)} \leftarrow
          \mathcal{R}_i^{(t)} \cdot a_i^{(t)} \cdot
          (1 + \lambda \cdot \widetilde{\mathrm{align}}_i^{(t)})$
         \hfill (Eq.~\ref{eq:score})
  \STATE $\mathcal{S}_t \leftarrow$ top-$B$ samples by $s_i^{(t)}$;
         fine-tune $\theta_{t-1} \to \theta_t$ on $\mathcal{S}_t$
\ENDFOR
\RETURN $\theta_T$
\end{algorithmic}
\end{algorithm}

\subsection{Stability of the trajectory anchor}
\label{sec:method-theory}

The anchor is extracted at each epoch from a random $1{,}024$-sample
probe; a natural concern is whether it is a meaningful structural
object or an artefact of that randomness. Theorem~\ref{thm:stability}
rules out the latter: the per-epoch anchor change is bounded by a
quantity that decays with the learning rate, and its cumulative
drift over training is controlled by the learning-rate schedule.

\begin{theorem}[Anchor stability]
\label{thm:stability}
Fix a layer $l$ and let $\Sigma^{(t)} := \Sigma_l^{(t)}$. Assume
\emph{(1)} bounded covariance change,
$\|\Sigma^{(t+1)} - \Sigma^{(t)}\|_F \le C_\Sigma\,\eta_t$;
\emph{(2)} a positive eigenvalue gap,
$\gamma^{(t)} := \lambda_1^{(t)} - \lambda_2^{(t)} \ge
\gamma_{\min} > 0$; and \emph{(3)} sign calibration,
$\langle v_l^{(t)}, \bar{\Delta h}_l^{(t)} \rangle > 0$ with
$\bar{\Delta h}_l^{(t)}$ varying slowly between epochs. Then the
sign-calibrated trajectory anchor satisfies the per-epoch bound
\begin{equation}
  \boxed{\;\bigl\| v_l^{(t+1)} - v_l^{(t)} \bigr\|
  \;\le\;
  \frac{2\, C_\Sigma}{\gamma_{\min}}\,\eta_t\;}
  \label{eq:stability-bound}
\end{equation}
and, by summation, the cumulative bound
$\|v_l^{(t+m)} - v_l^{(t)}\| \le
(2C_\Sigma/\gamma_{\min})\sum_{s=t}^{t+m-1}\eta_s$. Under any
decaying schedule the learning-rate tail $\sum_{s\ge t}\eta_s$
shrinks with $t$, so the anchor is $\varepsilon$-stationary: for
every $\varepsilon>0$ there is an epoch beyond which all subsequent
anchors lie within $\varepsilon$ of one another.
\end{theorem}

The proof applies the Davis--Kahan $\sin\Theta$
theorem~\cite{daviskahan} to $\Sigma^{(t)}, \Sigma^{(t+1)}$ and
converts the resulting angular bound to a Euclidean one via sign
calibration; the full argument, the assumption interpretations, and
a corollary transferring the bound to the selection score
(Cor.~\ref{cor:score-stability}) are given in
App.~\ref{app:proof-stability}. The bound concerns the anchor
together with a generic nonnegative base ranking score, so the
stability guarantee it yields for the \textsc{TADS} criterion is not
tied to the particular composite reward, advantage estimate, or
reward--advantage consensus rule used in our implementation. We verify
the theorem empirically in Sec.~\ref{sec:thm-empirical}, measuring
the per-epoch anchor change and the eigenvalue gap across training;
its broader consequences for reproducibility and runtime diagnostics
are discussed in Sec.~\ref{sec:discussion}.
